\chapter{봉투에 넣은 예측}

11장의 비교에는 작은 약점이 있어요.
점수도 검사도 같은 여섯 달의 기록에서 나왔다는 거예요.
허락 지도로 만든 점수가 허락을 많이 받은 지갑과 닮은 건 어찌 보면 당연해요.
이미 본 답안지로 채점하는 것과 조금 비슷하지요.

더 어려운 시험이 있어요.
\textbf{미래를 맞히는} 시험이에요.

\section{봉투 실험}

날씨 예보관을 떠올려 보세요.
“내일 비가 온다”고 말한 예보관이 정말 실력이 있는지는 어떻게 알까요?
오늘 예보를 봉투에 넣어 봉해 두고, 내일 진짜 날씨와 비교하면 돼요.
내일 날씨를 보고 나서 예보를 고칠 수 없게 하는 것이 중요해요.

저도 똑같이 했어요.

\begin{enumerate}
\item \textbf{2026년 2월 28일 밤 11시 59분 59초}까지의 기록만으로 지도를 그리고 점수를 매겨요.
      이 점수가 봉투에 든 예보예요. 이날을 \textbf{동결일}이라고 해요.
\item 그다음 \textbf{3월부터 5월까지} 석 달 동안 실제로 일어난 일을 세어요.
\item 예보 줄과 실제 줄을 켄달의 $\tau$로 비교해요.
\end{enumerate}

10장의 달력에서 “13장”이라고 적힌 줄이 바로 이 실험이에요.
봄에 채점했으니 \textbf{봄 홀드아웃}이라고 불러요.
“홀드아웃”은 채점에 쓸 기록을 따로 떼어 놓았다는 뜻이에요.

\section{무엇을 맞히나요?}

이번에 줄을 세운 것은 GMX 거래자들이 아니라, 동결일에 누군가에게서 살아 있는 허락을 받아 둔 통들이에요.
이 통들을 \textbf{스펜더 코호트}라고 불러요.
모두 1,335개였고, 그중 1,174개가 로봇(앱), 161개가 사람의 통이었어요.

맞힐 것은 두 가지예요.

\begin{itemize}
\item \textcolor{allowcol}{\textbf{새 허락}}: 3\textasciitilde{}5월에 새로 생긴 “허락한 통 → 이 통” 쌍이 몇 개인가.
      새 친구가 몇 명이나 문을 열어 주었는가예요. 1,335개 가운데 222개가 새 허락을 받았어요.
\item \textcolor{sendcol}{\textbf{새 송금자}}: 3\textasciitilde{}5월에 처음으로 이 통에 송금한 통이 몇 개인가.
      175개가 새 송금자를 맞았어요.
\end{itemize}

\section{결과}

이번에는 이 책에 나온 점수를 모두 한 줄에 세웠어요.
쪽지 놀이로 만든 점수 여섯 개와, 4장에서 만든 가장 단순한 점수 두 개예요.
가장 단순한 점수는 동결일까지 받은 화살표를 \emph{세기만} 한 거예요.
받은 허락의 수와 받은 송금의 수, 두 가지가 있어요.

\begin{figure}[h]
\centering
\begin{tikzpicture}
\begin{axis}[
  xbar, bar shift=0pt, bar width=6.5pt, width=0.86\linewidth, height=5.0cm,
  xmin=0, xmax=0.86, xtick={0,0.2,0.4,0.6,0.8},
  enlarge y limits=0.08,
  nodes near coords, nodes near coords style={font=\tiny},
  every node near coord/.append style={/pgf/number format/fixed, /pgf/number format/fixed zerofill, /pgf/number format/precision=3},
  tick label style={font=\scriptsize}, title style={font=\small\bfseries},
  title={새 허락 맞히기},
  symbolic y coords={L0, S, D, A, C, E, L1, G}, ytick={L0, S, D, A, C, E, L1, G},
  yticklabels={뒷면 랭크, 씨앗 랭크, 받은 송금 수, 송금 랭크, 동전 랭크, 엔도스랭크, 앞면 랭크, 받은 허락 수},
]
\addplot[fill=greycol!55, draw=greycol] coordinates {(0.711,G) (0.117,D)};
\addplot[fill=allowcol, draw=allowcol] coordinates {(0.479,L1) (0.394,E)};
\addplot[fill=leafline!70, draw=leafline] coordinates {(0.295,C)};
\addplot[fill=sendcol, draw=sendcol] coordinates {(0.123,A) (0.114,S) (0.044,L0)};
\end{axis}
\end{tikzpicture}

\medskip
\begin{tikzpicture}
\begin{axis}[
  xbar, bar shift=0pt, bar width=6.5pt, width=0.86\linewidth, height=5.0cm,
  xmin=0, xmax=0.86, xtick={0,0.2,0.4,0.6,0.8},
  enlarge y limits=0.08,
  nodes near coords, nodes near coords style={font=\tiny},
  every node near coord/.append style={/pgf/number format/fixed, /pgf/number format/fixed zerofill, /pgf/number format/precision=3},
  tick label style={font=\scriptsize}, title style={font=\small\bfseries},
  title={새 송금자 맞히기},
  symbolic y coords={S, L0, C, E, A, L1, G, D}, ytick={S, L0, C, E, A, L1, G, D},
  yticklabels={씨앗 랭크, 뒷면 랭크, 동전 랭크, 엔도스랭크, 송금 랭크, 앞면 랭크, 받은 허락 수, 받은 송금 수},
]
\addplot[fill=greycol!55, draw=greycol] coordinates {(0.472,D) (0.404,G)};
\addplot[fill=allowcol, draw=allowcol] coordinates {(0.372,L1) (0.283,E)};
\addplot[fill=leafline!70, draw=leafline] coordinates {(0.277,C)};
\addplot[fill=sendcol, draw=sendcol] coordinates {(0.344,A) (0.270,L0) (0.264,S)};
\end{axis}
\end{tikzpicture}
\caption{봄 홀드아웃(스펜더 1,335개). 막대가 길수록 미래를 잘 맞힌 거예요. 회색은 받은 화살표를 세기만 한 점수, \textcolor{allowcol}{파랑}은 허락 지도를 걷는 점수, \textcolor{sendcol}{주황}은 송금 지도를 걷는 점수, \textcolor{leafline}{초록}은 둘을 섞어 걷는 동전 랭크예요.}
\end{figure}

\subsection*{세기만 한 점수가 가장 잘 맞혔어요}

두 그래프 모두 맨 위에 회색 막대가 있어요.
새 허락 맞히기에서는 받은 허락의 수가 \textbf{0.711}로 가장 높았어요.
새 송금자 맞히기에서는 받은 송금의 수가 \textbf{0.472}로 가장 높았어요.

쪽지 놀이로 만든 점수는 어느 것도 세기만 한 점수를 넘지 못했어요.
허락 지도를 걷는 앞면 랭크(0.479)와 엔도스랭크(0.394)는 받은 허락의 수보다 한참 낮았어요.
송금 지도를 걷는 송금 랭크(0.344)도 받은 송금의 수보다 낮았어요.
세기와의 차이를 다시 뽑기로 재 보니 95\% 범위가 모두 0보다 한참 아래였어요.
우연이 아니라는 뜻이에요.

왜 그럴까요?
이미 허락을 많이 받은 앱은 앞으로도 허락을 많이 받아요.
사람들이 많이 쓰는 앱을 새 사람들도 쓰기 때문이에요.
“부자가 더 부자가 된다”는 말처럼요.
받은 화살표를 세면 이것을 바로 알 수 있어요.
쪽지 놀이는 한 친구의 믿음을 그 친구가 허락한 모든 통에 나누어 주는데, 이번 시험에서는 그렇게 나누는 것이 맞히는 데 도움이 되지 않았어요.

\subsection*{어느 지도를 걷느냐가 중요해요}

이번에는 막대의 색을 보세요.
새 허락 맞히기에서는 파란 막대가 위쪽에, 주황 막대가 아래쪽에 모여 있어요.
송금 기록만 쓰는 점수는 아무리 잘해도 0.123이었어요.
송금 기록만으로는 누가 새 허락을 받을지 거의 알 수 없다는 뜻이에요.
두 지도를 섞어 걷는 동전 랭크는 그 사이에 있었어요.

새 송금자 맞히기는 조금 달라요.
받은 송금의 수가 맨 위에 있지만, 그다음 자리는 받은 허락의 수와 앞면 랭크가 차지했어요.
허락을 많이 받는 앱은 대개 많은 사람이 쓰는 앱이라서, 처음 송금하는 사람도 많이 맞기 때문이에요.

그래서 이 시험으로 “엔도스랭크와 송금 랭크 가운데 누가 더 좋은가”를 가리지는 않아요.
두 점수는 다른 지도를 걷고 다른 것을 재요.
새 허락은 허락 화살표가 새로 생긴 것이니, 허락 지도를 걷는 점수가 잘 맞히는 게 자연스러워요.
두 점수를 맞대어 본 숫자는 책 끝의 「엔도스랭크와 송금 랭크를 맞대어 보면」(\pageref{ch:er-awp-side}쪽)에 따로 모아 두었어요.

이 비교들은 대부분 미리 정해 둔 것이 아니에요.
엔도스랭크와 송금 랭크 점수는 봄 채점 결과가 나온 뒤에 매겼거든요.
그래서 저는 숫자를 범위와 함께 그대로 알리기만 하고, 이것으로 무엇을 증명했다고 말하지는 않아요.
앞면 랭크와 받은 허락 수의 비교만은 9월 14일에 미리 정해 두었어요.

\subsection*{송금 랭크가 말할 수 없는 통}

8장에서 엔도스랭크는 허락 지도에 없는 지갑에 대해 아무 말도 할 수 없다고 했지요.
송금 랭크도 마찬가지예요.
스펜더 1,335개 가운데 318개는 동결일까지 다른 통과 송금을 주고받은 기록이 없어서 송금 랭크 점수가 0이었어요.
그런데 이 318개 가운데 22\%가 봄에 새 허락을 받았어요.
나머지 통들은 15\%였는데 말이에요.
송금 랭크는 새 허락을 받을 통이 많은 무리를 줄의 맨 뒤에 세운 셈이에요.

그렇다면 송금 랭크가 받은 송금의 수보다 낮은 것도 이 0점 때문일까요?
그래서 동결일 전에 송금을 받아 본 977개만 남기고 다시 재 보았어요.
그래도 새 송금자 맞히기에서 송금 랭크는 받은 송금의 수보다 한참 낮았어요(0.388 대 0.701).
엔도스랭크도 새 허락 맞히기에서 받은 허락의 수보다 낮았고요(0.340 대 0.679).
쪽지 놀이가 세기를 넘지 못한 것은 말할 수 없는 통 때문이 아니었어요.
이 점검도 결과를 본 뒤에 한 것이에요.

\section{동전 랭크는 약속을 지켰을까?}

9장에서 9월 14일에 적어 둔 약속을 기억하나요?
동전 랭크는 새 허락 맞히기에서 앞면 랭크보다 못하지 않고, 새 송금자 맞히기에서 뒷면 랭크보다 나아야 했어요.

\begin{itemize}
\item 새 허락: 동전 랭크 0.295, 앞면 랭크 0.479. 차이가 $-0.184$로 한참 못했어요. \textcolor{roseline}{\textbf{×}}
\item 새 송금자: 동전 랭크 0.277, 뒷면 랭크 0.270. 차이가 $+0.007$인데 95\% 범위가 $-0.012$에서 $0.028$로 0을 넘나들어요. 낫다고 할 수 없어요. \textcolor{roseline}{\textbf{×}}
\end{itemize}

두 조건 모두 넘지 못했어요.
동전 랭크는 9월 14일의 약속을 지키지 못했어요.
$\lambda$를 0.25나 0.75로 바꿔도 마찬가지였어요.

\section{그런데 눈에 띈 것이 하나 있었어요}

결과를 들여다보다가 약속에 없던 것이 눈에 띄었어요.
동전 랭크를 뒷면 랭크하고만 비교하면, 새 허락 맞히기에서 동전 랭크가 0.295로 뒷면 랭크의 0.044보다 \textbf{0.251}이나 높았어요.
새 송금자 맞히기에서는 뒷면 랭크와 비슷했고요.

“송금 지도에 허락 화살표를 섞으면 송금 지도만 걷는 것보다 나아진다.”
이것은 사실 제 논문 제목이 말하고 싶은 바로 그 이야기예요.
하지만 저는 이것을 결과라고 부를 수 없었어요.
\emph{결과를 보고 나서 찾은 것}이기 때문이에요.
송금 랭크와 견주어 보면 더 조심해야 해요.
새 허락에서는 동전 랭크가 앞섰지만(0.295 대 0.123), 새 송금자에서는 뒤졌거든요(0.277 대 0.344).
왜 결과를 보고 나서 찾은 것이 문제인지, 그래서 제가 무엇을 했는지가 다음 장의 이야기예요.

\begin{deeper}[진짜 숫자]
\begin{center}
\small
\begin{tabular}{@{}lcc@{}}
\toprule
점수 (2월 28일 동결) & 새 허락 $\tau$ & 새 송금자 $\tau$ \\
\midrule
엔도스랭크 & 0.394 & 0.283 \\
송금 랭크 & 0.123 & 0.344 \\
앞면 랭크 (동전 랭크, $\lambda=1$) & 0.479 & 0.372 \\
뒷면 랭크 (동전 랭크, $\lambda=0$) & 0.044 & 0.270 \\
동전 랭크 ($\lambda=0.5$) & 0.295 & 0.277 \\
씨앗 랭크 & 0.114 & 0.264 \\
받은 허락 수 & 0.711 & 0.404 \\
받은 송금 수 & 0.117 & 0.472 \\
\bottomrule
\end{tabular}
\end{center}
세기와의 차이는 엔도스랭크 $-$ 받은 허락 수가 새 허락 $-0.317$ [$-0.354$, $-0.283$], 송금 랭크 $-$ 받은 송금 수가 새 송금자 $-0.127$ [$-0.150$, $-0.099$]예요(둘 다 결과를 본 뒤 계산).
앞면 랭크 $-$ 받은 허락 수는 새 허락 $-0.233$ [$-0.291$, $-0.183$]으로, 9월 14일에 미리 정해 둔 비교예요.
송금을 받아 본 스펜더 977개에서는 송금 랭크 $-$ 받은 송금 수가 새 송금자 $-0.314$ [$-0.353$, $-0.276$]였어요.
허락한 양을 그대로 쓰는 앞면 랭크는 봄에 두 시험 모두 엔도스랭크보다 높았어요.
스펜더는 블록체인 노드에 코드가 있는지 물어서 계약 1,174개, 사람의 지갑 161개(계약 기능을 덧붙인 새 방식의 지갑 1개 포함)로 나누었어요.
\end{deeper}

\begin{learned}
\begin{itemize}
\item 2월 28일에 점수를 봉해 두고, 3\textasciitilde{}5월에 실제로 일어난 일로 채점했어요.
\item 받은 화살표를 세기만 한 점수가 가장 잘 맞혔어요. 쪽지 놀이로 만든 점수는 어느 것도 세기를 넘지 못했어요.
\item 새 허락은 허락 지도를 걷는 점수가 잘 맞혔고, 송금 기록만 쓰는 점수는 거의 맞히지 못했어요.
\item 송금 기록이 없는 통을 빼고 다시 재도, 송금 랭크는 받은 송금의 수보다 한참 낮았어요.
\item 동전 랭크는 9월 14일의 약속을 지키지 못했어요.
\item 동전 랭크가 뒷면 랭크보다 새 허락을 잘 맞힌 것은 결과를 보고 나서 찾은 것이라, 아직 결과라고 부를 수 없었어요.
\end{itemize}
\end{learned}
